\chapter{CRONOGRAMA}
\label{cap:cronograma}

O cronograma apresentado no Quadro~\ref{qua:cronograma} organiza as etapas de 
execução deste trabalho ao longo dos semestres de 2026.2 
(Projeto Supervisionado I) e 2027.1 (Projeto Supervisionado II).

\begin{table}[htb]
\caption{Cronograma de execução do trabalho}
\label{qua:cronograma}

\renewcommand{\arraystretch}{1.4}
\begin{tabularx}{\textwidth}{|X|c|c|c|c|c|c|c|c|c|c|}
  \hline
  \multirow{2}{*}{\textbf{Atividade}} &
  \multicolumn{5}{c|}{\textbf{2026.2 (PS I)}} &
  \multicolumn{5}{c|}{\textbf{2027.1 (PS II)}} \\
  \cline{2-11}
  & \textbf{Ago} & \textbf{Set} & \textbf{Out} & \textbf{Nov} & \textbf{Dez}
  & \textbf{Fev} & \textbf{Mar} & \textbf{Abr} & \textbf{Mai} & \textbf{Jun} \\
  \hline

  Definição do tema e problema de pesquisa
  & & $\checkmark$ & $\checkmark$  & & & & & & & \\
  \hline

  Revisão bibliográfica e referencial teórico
  & & $\checkmark$ & $\checkmark$ & $\checkmark$ & & & & & & \\
  \hline

  Redação da metodologia
  & & & $\checkmark$ & $\checkmark$  & & & & & & \\
  \hline

  Escrita e formatação do Projeto de Pesquisa
  & & & $\checkmark$ & $\checkmark$ & & & & & & \\
  \hline

  Entrega e defesa do Projeto de Pesquisa (PS I)
  & & & & & $\checkmark$ & & & & & \\
  \hline

  Desenvolvimento e calibração do protótipo
  & & & & & & $\checkmark$ & & & & \\
  \hline

  Implementação dos três geradores
  & & & & & & $\checkmark$ & $\checkmark$ & & & \\
  \hline

  Implementação e validação da BFS
  & & & & & & & $\checkmark$ & $\checkmark$ & & \\
  \hline

  Preparação ética e submissão ao CEP
  & & & & & & & & $\checkmark$ & $\checkmark$ & \\
  \hline

  Testes piloto e ajustes dos parâmetros
  & & & & & & & & $\checkmark$ & & \\
  \hline

  Experimentos computacionais
  & & & & & & & & $\checkmark$ & $\checkmark$ & \\
  \hline

  Avaliação com participantes
  & & & & & & & & & $\checkmark$ & \\
  \hline

  Análise e interpretação dos resultados
  & & & & & & & & & $\checkmark$ & $\checkmark$ \\
  \hline

  Redação do TCC completo
  & & & & & & & & $\checkmark$ & $\checkmark$ & \\
  \hline

  Revisão e formatação final (ABNT)
  & & & & & & & & & $\checkmark$ & \\
  \hline

  Entrega e defesa do TCC (PS II)
  & & & & & & & & & & $\checkmark$ \\
  \hline

\end{tabularx}

\fonte{Elaborado pelo autor (\oano).}
\end{table}

As etapas previstas no cronograma são descritas a seguir:

\begin{description}

  \item[\textbf{Definição do tema, problema e objetivos:}]
  Delimitação do objeto de estudo, formulação do problema de pesquisa e 
  estabelecimento dos objetivos geral e específicos.

  \item[\textbf{Revisão bibliográfica e referencial teórico:}]
  Levantamento e análise de publicações sobre geração procedural de conteúdo, 
  jogos de plataforma 2D, algoritmos genéticos e avaliação de fases geradas 
  automaticamente.

  \item[\textbf{Elaboração da metodologia:}]
  Definição das estratégias de geração, da representação computacional das 
  fases, dos procedimentos experimentais e das métricas de avaliação.

  \item[\textbf{Redação e apresentação do PS I:}]
  Organização, revisão e apresentação do projeto de pesquisa, contemplando sua
  fundamentação teórica e metodologia proposta.

  \item[\textbf{Desenvolvimento e calibração do protótipo:}]
  Implementação do ambiente de jogo de plataforma 2D, incluindo movimentação,
  saltos, colisões e representação das fases. Calibração preliminar das 
  características físicas do personagem.

  \item[\textbf{Implementação dos três geradores:}]
  Desenvolvimento das estratégias de geração aleatória, geração baseada em 
  regras e geração por algoritmo genético, utilizando uma representação comum 
  das fases.

  \item[\textbf{Implementação e validação da BFS:}]
  Desenvolvimento do procedimento automatizado de verificação de 
  completabilidade por busca em largura, com testes de funcionamento e 
  reprodução das soluções encontradas.

  \item[\textbf{Preparação ética e submissão ao CEP:}]
  Elaboração dos instrumentos de coleta, do Termo de Consentimento Livre e
  Esclarecido e encaminhamento da pesquisa para apreciação ética, conforme as
  exigências institucionais aplicáveis.

  \item[\textbf{Testes piloto e ajustes dos parâmetros:}]
  Execução de experimentos preliminares para verificar o funcionamento dos 
  geradores, os custos computacionais e a adequação dos procedimentos de 
  avaliação.

  \item[\textbf{Experimentos computacionais:}]
  Realização das execuções independentes previstas para cada estratégia, com 
  registro de completabilidade, diversidade estrutural e desempenho 
  computacional.

  \item[\textbf{Avaliação com participantes:}]
  Realização das sessões de jogo após as autorizações aplicáveis, com coleta de
  dados de desempenho e percepções dos participantes por meio de questionários.

  \item[\textbf{Análise e interpretação dos resultados:}]
  Organização e análise dos dados quantitativos e qualitativos, visando 
  identificar diferenças entre as estratégias investigadas e discutir suas 
  características e limitações.

  \item[\textbf{Redação, revisão e defesa final:}]
  Consolidação dos resultados, redação das discussões e conclusões, revisão do 
  documento e apresentação do Trabalho de Conclusão de Curso.

\end{description}
